% from main.tex
%%%%%%%%%%  scan idx 406  |  folio 393  %%%%%%%%%%
Soit \(E\) un \(A\)-Module localement libre de rang fini \(n\). Il existe un torseur \(P\) sous \(\GL(n,A)\) tel que \(E=P\times_{\GL(n,A)} (A)^n\), où \((A)^n\) est considéré comme \(\GL(n,A)\)-Module par la représentation standard. On peut alors définir un \(\lambda\)-homomorphisme de \(R_{\mathbb Z}(\GL(n))\) dans \(K\), en associant à toute représentation \(\rho\) de \(\GL(n)\) dans un \(\mathbb Z\)-module libre de rang \(p\) la classe du \(A\)-Module localement libre
\(P\times_{\GL(n,A)} (A)^p\), où \((A)^p\) est considéré comme \(\GL(n,A)\)-module par la représentation \(\rho\). Dans cet homomorphisme, la classe de \(E\) est donc l'image de la représentation standard. Pour vérifier une relation polynomiale entre les \(\lambda^i(E)\), il suffit donc que cette relation soit vérifiée par la représentation standard dans \(R_{\mathbb Z}(\GL(n))\).

On peut faire un raisonnement analogue avec deux \(A\)-Modules localement libres \(E,F\) de rangs \(n,p\) ; on peut alors trouver un torseur \(P\) sous \(\GL(n,A)\times\GL(p,A)\) tel que \(E=P\times_{\GL(n,A)\times\GL(p,A)} (A)^n\), où \((A)^n\) est considéré comme \(\GL(n,A)\times\GL(p,A)\)-Module en faisant opérer \(\GL(n,A)\) par la représentation standard et \(\GL(p,A)\) par la représentation triviale, et réciproquement pour \(F\). On définit ensuite un \(\lambda\)-homomorphisme de \(R_{\mathbb Z}(\GL(n)\times\GL(p))\) dans \(K\) de la même façon que plus haut, et on se ramène ainsi à vérifier les relations entre les \(\lambda^i(E)\), \(\lambda^j(F)\) pour les éléments correspondants de \(R_{\mathbb Z}(\GL(n)\times\GL(p))\). On est donc ramené à montrer que pour tout schéma en groupes \(G\) sur \(\Spec\mathbb Z\), produit fini de schémas en groupes linéaires, le pré-\(\lambda\)-anneau \(R_{\mathbb Z}(G)\) est un \(\lambda\)-anneau. Pour ce faire, on peut référer à [5], où \(R_{\mathbb Z}(G)\) est calculé explicitement. Si l'on ne veut pas utiliser ce résultat, on peut raisonner de façon analogue à 3.2, en utilisant 1.14.

